\end{center}

\subsection{The measured cell budget}\label{sec:budget}

\begin{plain}
Below is what the machine actually costs, measured by running the industry-standard
open-source tools --- Yosys to translate the design into logic gates, nextpnr to place
those gates on a real chip --- targeting a Lattice iCE40 UP5K, a chip you can buy on a
breakout board for about \$25.

The entire arithmetic core of this computer is \emph{seventy-two} logic cells, out of$5{,}280$ available. It uses $1.4\%$ of a \$25 chip.
\end{plain}

\begin{spec}[Measured; iCE40 UP5K SG48, Yosys 0.67 $+$ nextpnr]
\begin{center}\small
\begin{tabular}{@{}lrrl@{}}
\toprule
\textbf{Build} & \textbf{Logic cells} & \textbf{$F_{\max}$ (MHz)} & \textbf{Contains}\\
\midrule
$F$ only                 & 16 & 230.84 & one Fourier gate \\
$F_p + F_f$              & 19 & 230.84 & both Fourier gates \\
$F$ and $S$ (4 opcodes)  & 40 &  86.60 & the single-qutrit Cliffords \\
$\mathrm{CX}$ only       & 33 &  72.40 & the entangler, both directions \\
$\sigma^5{=}Z$ only      & 20 & 144.07 & the translation \\
all Clifford, no $Z$     & 65 &  60.80 & \textcolor{bad}{cannot leave the origin} \\
\textbf{full frame unit} & \textbf{72} & \textbf{60.80} & \textcolor{good}{the machine} \\
\midrule
\multicolumn{4}{@{}l}{\itshape For comparison, elsewhere in the design:}\\
$\mathrm{CX}$, bare loadable        &  23 & 72.40 & \cmd{w33\_pass2752\_cx\_loadable\_frame.sv}\\
$36$-lane mixer, serial            & 4048 & 19.65 & \cmd{w33\_pass2612\_serial\_mixer.sv}\\
$36$-lane mixer, parallel          & \multicolumn{2}{l}{13{,}965 LUT4 $+$ 799 carry} &
    \textcolor{bad}{does not fit any iCE40}\\
\bottomrule
\end{tabular}
\end{center}
\end{spec}

\begin{plain}
The last two rows are worth a moment. The ``mixer'' is the part that combines $36$
signals at once. Built the obvious way --- all $36$ lanes in parallel --- it needs
about $14{,}000$ logic cells, nearly twice what the largest chip in this family has. Built
serially, one lane at a time, it needs $4{,}048$: about a third of the area, at
$1/36$ of the speed. That is the kind of trade a blueprint exists to make explicit.
\end{plain}

\begin{gotwrong}
\textbf{Every cell count in this project before Pass 2753 was wrong, and the reason is
instructive.} We reported $13$ cells for $F$ and $S$, $8$ for $\mathrm{CX}$, $42$ for the
ISA. Those modules had no load port. By the reachability theorem of \S\ref{sec:isa}, a
register driven only by Clifford opcodes can never leave $(0,0,0,0)$ --- so the synthesis
tool \emph{proved the state was constant and deleted six of the eight flip-flops}. We were
measuring the area of nothing.

\medskip\noindent
Two independent teams shipped this same defect within four hours of each other. Both had
\emph{exhaustive} testbenches --- one covered all $81^2$ pairs of frames --- and both
passed, because the testbenches drove the \emph{combinational} logic and never
instantiated the sequential wrapper.

\medskip\noindent
\textbf{Simulation cannot detect this class of defect at all.} A module that folds to a
constant still simulates correctly. It just does nothing. Only a synthesis frontend asks
whether the state can move.

\medskip\noindent
One figure survived unchanged: $\sigma^5=Z$ was reported at $21$ cells and re-measures at
$20$. It was always real --- because it is the one instruction that is not linear, and so
the one instruction that could always move the state. The theory predicted exactly which
number would survive.
\end{gotwrong}

% =====================================================================================
\section{The physical layer}
% =====================================================================================

\begin{plain}
Everything above is logic --- it would work on any hardware that can implement the eight
operations. This section is about doing it with light, which is the intended
implementation, and about what has already been demonstrated in a laboratory as opposed
to designed on paper.

A single photon carries several independent properties. Two of them are useful here:
\emph{when} it arrives (its time bin) and \emph{what colour} it is (its frequency bin).
Chop time into three slots and colour into three bands, and one photon carries two
qutrits --- the ``past'' and ``future'' registers of \S\ref{sec:physics}.
\end{plain}

\begin{center}
\begin{tikzpicture}[font=\footnotesize]
  \draw[->,line width=0.8pt] (0,0) -- (9.0,0) node[right,font=\scriptsize] {time};
  \draw[->,line width=0.8pt] (0,0) -- (0,3.2) node[above,font=\scriptsize] {frequency};
  \foreach \i in {0,1,2}{
    \draw[dashed,draw=rule] (0,\i*0.95+0.5) -- (8.6,\i*0.95+0.5);
    \draw[dashed,draw=rule] (\i*2.7+0.9,0) -- (\i*2.7+0.9,3.0);
  }
  \foreach \i in {0,1,2}{
    \node[font=\scriptsize,text=future,left] at (-0.05,\i*0.95+0.95) {$f_\i$};
    \node[font=\scriptsize,text=past,below] at (\i*2.7+1.35,-0.08) {$t_\i$};
  }
  \foreach \i in {0,1,2}{
    \fill[past!60,rounded corners=1pt]
      (\i*2.7+1.0,\i*0.95+0.62) rectangle (\i*2.7+1.7,\i*0.95+1.28);
  }
  \node[align=left,text width=5.0cm,font=\scriptsize] at (12.3,1.7)
   {\textbf{The \textsc{sum} gate}\\[3pt]
    $\ket{f,t}\mapsto\ket{f,\;t+f \bmod 3}$\\[4pt]
    A chirped fibre Bragg grating delays each colour by a different amount, so the
    photon's \emph{colour} controls how far its \emph{arrival time} shifts. That is a
    controlled-add, implemented with one passive optical element.};
  \node[align=left,text width=5.0cm,font=\scriptsize,text=rule] at (12.3,-0.5)
   {Demonstrated: Imany \emph{et al.}, \emph{npj Quantum Information} \textbf{5}, 59
    (2019). Three $3$\,ns time bins at $6$\,ns spacing, three frequency bins at
    $380$\,GHz, $-2$\,ns/nm dispersion, $18$\,ns wraparound. Measured basis fidelity
    $0.92\pm0.01$; certified entanglement of formation $\ge 1.19\pm0.12$ ebits.};
\end{tikzpicture}
\end{center}

\begin{spec}[Evidence boundary --- what is measured and what is not]
\textbf{Measured in a laboratory:} the deterministic single-photon qutrit \textsc{sum}
gate, at fidelity $0.92\pm0.01$, with the entangled output certified.

\smallskip\noindent
\textbf{Proved exactly:} the group theory, the geometry, the class atlas, the ISA
semantics, the reachability, the synthesis and placement figures.

\smallskip\noindent
\textbf{Neither, and stated as such:} source efficiency, insertion-loss engineering,
fault-tolerance thresholds, and the magic-state injection route. The $M_{36}$ resources
are typed \cmd{M36\_Q4\_RAW} in the controller and injection is \emph{refused} in RTL
until a proved ququart protocol supplies a validity assertion. That refusal is a designfeature: it makes the gap impossible to paper over.
\end{spec}

% =====================================================================================
\section{Why the exponent nine}
% =====================================================================================

\begin{plain}
A small but pretty result that shows how tightly the layers fit.

To tell which operation the machine just performed, you measure a quantity and look it up
in a table. The obvious quantity --- the ``trace'' --- has a flaw: it changes if the whole
state picks up an irrelevant overall phase, which physically means nothing. The fix used
in this project is to raise the trace to the \emph{ninth} power and divide by a
determinant. The ninth power was found empirically. Where does nine come from?
\end{plain}

\begin{spec}[Derived, Pass 2779]
For a $d$-dimensional representation, under $U\mapsto\lambda U$ we have
$\mathrm{Tr}(U^k)^e\mapsto\lambda^{ke}\,\mathrm{Tr}(U^k)^e$ and
$\det(U^k)\mapsto\lambda^{kd}\det(U^k)$, so
\[
  \Theta_k(g)=\frac{\mathrm{Tr}(U_g^k)^e}{\det(U_g^k)}
\]
is phase-invariant for all $\lambda$ \emph{exactly} when $e=d$. Here $d=9$, the dimension
of the two-qutrit state space --- which by \S\ref{sec:physics} is
$\dim\mathrm{End}(\mathbb{C}^3)$. Verified numerically: $e\in\{3,8,10\}$ fail, $e=9$ holds.

\smallskip\noindent
If the phase ambiguity were a \emph{finite} group $\mu_m$ one could use any
$e\equiv 9\pmod m$, possibly smaller. We computed the scalar subgroup of the generated
matrix group by collision sampling and it is $\mu_{12}$. Since $9\not\equiv 0\pmod{12}$,
the determinant is genuinely required, and $9$ is the \emph{smallest} exponent that
works: the sensor is minimal.
\end{spec}

\headline{$\mu_{12}=\mu_4\times\mu_3$. The $\mu_3$ is $\omega$, the qutrit alphabet.
The $\mu_4$ is the quadratic Gauss sum $\sum_j\omega^{j^2}=i\sqrt3$, which is
\emph{imaginary} precisely because $3\equiv 3 \pmod 4$ --- the same congruence that makes
time reversal a physical rather than a gauge symmetry.}

\begin{plain}
That last box is the kind of thing that makes this project worth doing. A practical
question --- ``why is the exponent nine in the measurement circuit?'' --- turns out to
have the same answer as an abstract one --- ``why does this geometry distinguish past
from future?''. Both are the statement that $-1$ is not a perfect square when you count
modulo $3$.
\end{plain}

\subsection{And for wider registers, it is usually three}

\begin{spec}[Pass 2791]
The phase group is not an artefact of the two-qutrit case. At $n=1$ the Clifford group is
small enough to enumerate exactly rather than sample:
\[
  |\langle F,S,X\rangle| = 2592 = \underbrace{12}_{\mu_{12}}\times
  \underbrace{216}_{|\mathrm{ASp}(2,3)|}\quad\text{--- consistent, so }\mu_{12}\text{ again.}
\]
For $n$ qutrits $d=3^n$, so the minimal exponent is the least $e\equiv 3^n \pmod{12}$, and
$3^n \bmod 12$ cycles $3,9,3,9,\dots$:
\[
\begin{array}{c|cccc}
n & 1 & 2 & 3 & 4\\\hline
d=3^n & 3 & 9 & 27 & 81\\
\text{minimal }e & \mathbf{3} & 9 & \mathbf{3} & 9
\end{array}
\]
Verified directly on $3\times3$ unitaries: $e=3$ is phase-invariant, $e\in\{2,4,9\}$ is not.
\end{spec}

\headline{Odd register widths need only a \emph{cube} of the trace, not a ninth power ---
a strictly cheaper measurement, alternating with width.}

% =====================================================================================
\section{The magic states, and the one thing the machine cannot yet do}\label{sec:magic}
% =====================================================================================

\begin{plain}
Everything described so far is, in a precise technical sense, \emph{not enough}. The
eight opcodes generate a large and useful group, but a machine restricted to those
operations can be simulated efficiently by an ordinary computer. To get a genuine quantum
advantage you need one more ingredient, traditionally supplied as a special prepared
state called a \emph{magic state}.

The substrate hands us $36$ candidates for free --- $36$ special directions in the
four-dimensional space, called the Witting rays. The open question, and the single
largest gap in this blueprint, is whether those $36$ can be \emph{purified}: whether many
noisy copies can be distilled into one clean one. Without that, the machine is a very
elegant classical simulator.
\end{plain}

\subsection{All thirty-six look identical from the inside}

% The optional argument is braced: an unbraced "]" inside $\mathbb{Z}[\omega]$ closes
% the optional argument early and LaTeX reports "Extra }, or forgotten $".
\begin{spec}[{Exact integer arithmetic in $\mathbb{Z}[\omega]$, Pass 2790}]
Only two overlap values occur among the $36$ rays:
\[
  |\langle i|j\rangle|^2 \in \{\,0,\ \tfrac13\,\},
\]
and every ray has the \emph{same} profile: $11$ orthogonal partners and $24$ partners at
overlap $1/3$. There is exactly one Gram profile among all $36$.

\smallskip\noindent
(The orthogonality graph is $11$-regular on $36$ vertices but \emph{not} strongly
regular: $\lambda\in\{1,2\}$, $\mu\in\{3,4\}$. Recorded because the near-miss invites a
false claim.)
\end{spec}

\begin{plain}
So no symmetry of the machine can tell one ray from another by looking at how they sit
